% SOURCE_AUTHORITY: sga5_fr_workpass.tex, lines 13642--13658 (LF-normalized unit).
% SOURCE_SHA256: 791F4EFFC5E02832D5D77ED03518C8156D6F07E4C8238B03545DB93D883FBB28
Podemos concluir agora a demonstração do lema 4.7. Primeiro, como o espaço vetorial $\fm/\fm^2$ é de dimensão 1, tem-se
\[
\Aut_k(\fm/\fm^2) = k^* \quad ,
\]
e, portanto, $\im(G \longrightarrow \Aut_k(\fm/\fm^2))$ é um subgrupo finito cíclico de $k^*$, gerado por uma raiz primitiva $c$-ésima da unidade. Denotemos por $\varepsilon$ uma raiz deste tipo. O subgrupo $\im(G \longrightarrow \Aut_k(\fm/\fm^2))$ é gerado, portanto, pelo automorfismo
\[
\chi : \fm/\fm^2 \longrightarrow \fm/\fm^2 \quad , \quad \xi \longmapsto \varepsilon\xi \quad .
\]
Como o elemento $\mu(f'_R)$ de $\bigotimes_k^{r-1}(\fm/\fm^2)$ é invariante por $G$, tem-se:
\[
\mu(f'_R) = \chi\mu(f'_R) = \chi\sum_i \xi_{1i} \otimes \xi_{2i} \otimes \cdots \otimes \xi_{r-1,i} = \sum_i \varepsilon\xi_{1i} \otimes \varepsilon\xi_{2i} \otimes \cdots \otimes \varepsilon\xi_{r-1,i} = \varepsilon^{r-1}\mu(f'_R) \quad ,
\]
Isto é,
\[
(1 - \varepsilon^{r-1})\mu(f'_R) = 0 \quad ,
\]
Como $\mu(f'_R)$ é diferente de zero (cond. I), resulta que $1 - \varepsilon^{r-1} = 0$, ou seja, $c \mid (r-1)$.
\end{prooftext}
